\documentclass[10pt,a4paper]{article}
\usepackage{amsmath,amssymb}
\usepackage[margin=22mm]{geometry}
\usepackage{enumitem}
\usepackage{booktabs,longtable,array}
\usepackage{xcolor}
\definecolor{linkblue}{RGB}{20,60,140}
\usepackage[colorlinks=true,linkcolor=linkblue]{hyperref}

\newcommand{\code}[1]{\texttt{#1}}
\newcommand{\mi}[1]{\boldsymbol{#1}}

\title{DNNF90 Specification\\[1mm]
{\large High-order partial-derivative propagation through neural networks\\
by multi-index partial Bell polynomials --- v1.4}}
\date{August 2026}
\author{}

\begin{document}
\maketitle

\tableofcontents
\bigskip

%=====================================================================
\section{Overview}

\subsection{Purpose and position}

DNNF90 is a Fortran 90/95 library that constructs the output of a
fully connected neural network together with its \textbf{mixed partial
derivatives of arbitrary order} in a single forward pass, and returns
the exact weight gradient of any loss built from those derivatives in a
single reverse pass. The core is the recursion for high-order
derivatives of composite functions organized by multi-index partial
Bell polynomials; no nested automatic differentiation and no finite
differences are used.

The software is released under the MIT license. It has two layers.

\begin{itemize}[leftmargin=1.4em]
\item \textbf{The library} (\code{lib/}): the embeddable mechanism.
      One-point evaluation, one-point seeded gradients, one optimizer
      step applied to a caller-owned array, committee estimation,
      Kalman-filter training, and the C bindings.
\item \textbf{The training application} (\code{app/}): the
      keyword-driven reference implementation. Minibatches,
      validation, early stopping, eleven optimizers, file I/O, MPI.
\end{itemize}

The boundary follows the separation of mechanism from policy
(Section~\ref{sec:boundary}).

\subsection{Scope}

\begin{itemize}[leftmargin=1.4em]
\item The hidden activation is $\tanh$ only, because the derivative
      tables are built from the Bell polynomials of $\tanh$.
\item The output layer is linear (identity) with a single output; the
      propagated derivatives are those of a linear read-out of the
      last hidden layer.
\item The input dimension $D_0$ and the maximum order $K$ are given in
      the input. The carried multi-index set is either dense (all
      $|\mi{\alpha}|\le K$) or the downward closure of a prescribed
      seed set.
\item The loss lives outside the library. The library receives
      $\partial L/\partial T_{\mi\alpha}$ (the \emph{seed} below) and
      returns $\partial L/\partial W$.
\end{itemize}

%=====================================================================
\section{Mathematical specification}

\subsection{Notation}

Layers $l=1,\dots,L$, width $n_l$ of layer $l$, weights
$W^{(l)}_{ji}$ ($i=0$ is the bias). Input
$\mi{x}\in\mathbb{R}^{D_0}$. Multi-index
$\mi\alpha=(\alpha_1,\dots,\alpha_{D_0})$,
$|\mi\alpha|=\sum_i\alpha_i$. The carried quantities are

\[
T^{(l,\mi\alpha)}_j \;=\; \frac{\partial^{|\mi\alpha|} z^{(l)}_j}{\partial x_1^{\alpha_1}\cdots \partial x_{D_0}^{\alpha_{D_0}}},
\qquad
S^{(l,\mi\alpha)}_j \;=\; \frac{\partial^{|\mi\alpha|} a^{(l)}_j}{\partial \mi{x}^{\mi\alpha}},
\]
where $a^{(l)}=W^{(l)}z^{(l-1)}$ and $z^{(l)}=\sigma(a^{(l)})$
($\sigma=\tanh$; the last layer is the identity).

\subsection{Forward propagation (high order)}
\label{sec:fwd}

The linear part is linear multi-index by multi-index:
\[
S^{(l,\mi\alpha)}_j \;=\; \sum_{i} W^{(l)}_{ji}\,T^{(l-1,\mi\alpha)}_i ,
\qquad
T^{(0)}_0 \equiv 1 \ (\text{the bias channel}).
\]
The nonlinear part is the high-order derivative of a composition,
expressed through the multi-index partial Bell polynomials
$B_{\mi\alpha,q}$:
\[
T^{(l,\mi\alpha)}_j \;=\; \sum_{q=1}^{|\mi\alpha|} \sigma^{(q)}\!\big(S^{(l,\mi 0)}_j\big)\;
B_{\mi\alpha,q}\big(S^{(l,\cdot)}_j\big),
\]
where $B_{\mi\alpha,q}$ is a polynomial in the lower multi-index
components of $S$, built by the recursion
$B_{\mi\alpha,q}=\sum_{\mi\beta} c_{\mi\alpha\mi\beta}\,S^{(\mi\beta)}B_{\mi\alpha-\mi\beta,q-1}$.
The coefficients $c$ and the referenced indices are tabulated once at
initialization. The input layer carries $T^{(1,\mi 0)}_j=x_j$,
$T^{(1,\mi e_j)}_j=1$, and zero otherwise.

\subsection{High-order derivatives of $\tanh$}

$\sigma^{(q)}$ closes as a polynomial in $t=\tanh$:
$\sigma^{(0)}=t$ and
$\sigma^{(q+1)}=(1-t^2)\,\mathrm{d}\sigma^{(q)}/\mathrm{d}t$,
carried as a recursion on coefficient arrays, so that
$q=0,\dots,K+1$ costs one polynomial evaluation per point.

\subsection{Reverse propagation (adjoint)}

Let the seed of a loss $L$ be
$\bar T^{(L,\mi\alpha)}_1=\partial L/\partial T^{(L,\mi\alpha)}_1$.
The adjoint quantities are
\[
\bar S^{(l,\mi\beta)}_j=\sum_{\mi\alpha} \bar T^{(l,\mi\alpha)}_j\, h^{(l)}_{j,\mi\alpha\mi\beta},
\qquad
\bar T^{(l-1,\mi\alpha)}_i=\sum_j W^{(l)}_{ji}\,\bar S^{(l,\mi\alpha)}_j,
\]
\[
\frac{\partial L}{\partial W^{(l)}_{ji}}=\sum_{\mi\alpha}\bar S^{(l,\mi\alpha)}_j\,T^{(l-1,\mi\alpha)}_i .
\]
$h$ is assembled from the $\sigma^{(q+1)}$ values saved in the forward
pass paired with the Bell tables. The gradient is accumulated in a
fused per-layer sweep as soon as $\bar S^{(l)}$ is complete, so
$\bar T,\bar S$ hold only one layer's plane at a time.

\subsection{Closure sets}

Given a seed multi-index set $\mathcal{A}$, the set actually carried is
its \textbf{downward closure}
$\{\mi\beta \mid \exists \mi\alpha\in\mathcal{A},\ \mi\beta\le\mi\alpha\}$.
A non-closed set does not close the recursion. The dense set
$|\mi\alpha|\le K$ is a special case of a closure. Indexing is
lexicographic order with binary search (\code{find\_alpha}), which
remains viable at fifty to a hundred descriptors.

%=====================================================================
\section{Software structure}

\subsection{The layer boundary: mechanism and policy}
\label{sec:boundary}

The library provides \textbf{mechanism}: evaluate one point, form one
point's weight gradient from a seed, apply one step to a caller-owned
array. The trainer holds \textbf{policy}: batch size, validation
cadence, stopping rules, file output.

The boundary is forced by the embedding cases. Density-functional and
molecular-dynamics host codes own their outer loop, their data flow
and their parallel decomposition, so a library that imposes an epoch
loop, a validation split or a checkpointing policy cannot be embedded.
Mechanism can stay pure and reentrant (one work space per thread is
the only requirement) and admits exact verification by finite
differences; no comparable check exists for a stopping rule. The
optimizer asymmetry (two rules in \code{lib/}, eleven in \code{app/})
follows from the same principle. The relationship mirrors that of
PyTorch proper to Lightning; \code{app/} ships as the library's
largest worked example.

\subsection{File layout}

\begin{longtable}{@{}p{0.30\textwidth}p{0.10\textwidth}p{0.54\textwidth}@{}}
\toprule
File & Lines & Role \\
\midrule
\endhead
\multicolumn{3}{@{}l}{\textbf{lib/} (the embeddable library)}\\
\code{multi\_index\_bell\_module} & 677 & multi-index generation, downward closure, Bell tables, high-order $\tanh$ derivatives, lexicographic indexing, table sets (\code{tabset\_t})\\
\code{net\_module}       & 322 & create/read/write/free of \code{net\_t}, evaluation work space, \code{net\_eval\_hod}\\
\code{train\_module}     & 400 & \code{net\_forward\_point} / \code{net\_backward\_point} / \code{net\_grad\_point}, the batch gradient \code{net\_grad\_batch}, gradient accumulators, \code{opt\_sgd\_step} / \code{opt\_adam\_step}\\
\code{committee\_module} & 105 & uncertainty estimation by a committee of networks\\
\code{kalman\_module}    & 239 & per-pattern extended-Kalman-filter training\\
\code{c\_api\_module} + \code{dnnf90.h} & 479 & the C bindings (22 functions)\\
\midrule
\multicolumn{3}{@{}l}{\textbf{app/} (the keyword-driven trainer)}\\
\code{global\_variables}  & --- & shared state\\
\code{io\_module}         & --- & input parsing, data reading, output formats\\
\code{lib\_net\_module}   & 210 & the lib network the trainer holds; weight synchronization and the entry points\\
\code{pinn\_module}       & 59  & residual evaluation and residual seeds (policy)\\
\code{init\_weight\_module} & --- & eight initializations\\
\code{optimizer\_module}  & --- & eleven optimizers\\
\code{sgd\_batch\_module} & --- & epoch loop, minibatches, early stopping, MPI averaging\\
\code{validation\_module} & --- & validation costs, RMSE/MAE, output\\
\code{hod\_check\_module} & --- & in-run finite-difference self-check\\
\code{api\_module}        & --- & Fortran API onto the trainer state\\
\bottomrule
\end{longtable}

\subsection{The synchronization invariant}

Because the trainer owns the master weight array, the invariant
``the library network must hold the current weights'' arises. It is
enforced mechanically, not by discipline: every write to \code{weight}
and \code{weight\_best} advances a generation counter, each
synchronization records the generation and the source array, and every
evaluation entry point compares the two before computing and stops on
a mismatch. The check is one integer comparison and stays enabled in
production builds, because a missed synchronization returns
\emph{plausible numbers computed from the previous weights} --- the
most dangerous failure mode this layer has.

\subsection{Uniqueness of the propagation engine}

There is \textbf{exactly one} propagation engine in the tree.
Training, the validation cost, the prediction output, the self-check
and the C API all pass through \code{net\_forward\_point} /
\code{net\_backward\_point} / \code{net\_grad\_batch}. The trainer
keeps its own weight array as the master copy (checkpoints, restarts,
MPI averaging and the optimizer states are expressed in it);
\code{lnet\_sync\_weights} copies it into the library network once per
batch iteration and at every point where the weights can have moved
(just before validation, at output writing, inside the
finite-difference loops of the self-check).

The entry points of \code{app/lib\_net\_module.f90}:

\begin{longtable}{@{}p{0.32\textwidth}p{0.62\textwidth}@{}}
\toprule
Entry & Meaning \\
\midrule
\endhead
\code{lnet\_sync\_weights(w)} & copy the trainer's weights into the library network\\
\code{lnet\_forward\_value(x,v)} & one point's value\\
\code{lnet\_forward\_hod(x,t)} & one point's value and all carried derivatives, into the caller's array \code{t}\\
\code{lnet\_nalpha()} & the number of carried multi-indices (so the caller can size arrays)\\
\code{lnet\_value\_grad(x,y,c,$\nabla$)} & accumulate one point's value-fit weight gradient\\
\code{lnet\_seed\_grad(seed,$\nabla$)} & accumulate the adjoint of an arbitrary seed\\
\code{lnet\_seed\_row(seed,row)} & write the adjoint of an arbitrary seed as one row (the natural-gradient metric)\\
\code{lnet\_value\_row(row)} & the row of a unit seed on the value component\\
\code{lnet\_batch\_grad(X,Y,c,$\nabla$)} & the value-fit gradient of a whole minibatch at once\\
\bottomrule
\end{longtable}

%=====================================================================
\section{Library API}
\label{sec:libapi}

\code{lib/} is the embeddable mechanism and is usable independently of
the training application; this section describes everything it
exposes. \code{make lib} builds the static library
(\code{build/libdnnf90.a}); \code{make shared} the shared one. From
Fortran, \code{.mod} files are compiler- and version-specific, so
embeddings that cross toolchains should use the C bindings
(Section~\ref{sec:capi}).

\subsection{The contract}

The training contract is per point:
\[
\text{\code{net\_grad\_point}}:\quad
\underbrace{\text{seed}_{\mi\alpha}=\frac{\partial L}{\partial T^{\mi\alpha}}}_{\text{assembled by the caller}}
\ \longmapsto\
\frac{\partial L}{\partial W}
\]
The content of the loss $L$, the chain rule from descriptors to
inputs, and the meaning of grids or collocation points all belong to
the caller; the library does not know them. There is no shared state,
and one work space per thread suffices for concurrent calls (the
library is built with \code{-frecursive}).

\subsection{\code{multi\_index\_bell\_module}: multi-indices and tables}

\begin{longtable}{@{}p{0.36\textwidth}p{0.58\textwidth}@{}}
\toprule
Element & Role \\
\midrule
\endhead
\code{init\_hod\_tables(D0,K,ns,seeds)} & build the shared tables. \code{ns=0}: the dense set ($|\mi\alpha|\le K$); \code{ns>0}: the downward closure of the seed set\\
\code{tabset\_t} & a table-set type; several networks can carry different multi-index sets at once\\
\code{tabset\_init}, \code{tabset\_from\_current}, \code{tabset\_free} & construct a table set (fresh, or copied from the shared tables) and free it\\
\code{alpha\_index(alpha)} & slot number of a multi-index (lexicographic order + binary search)\\
\code{alpha\_list}, \code{alpha\_deg}, \code{ind\_e1} & the list of multi-indices, their orders, the slots of the unit multi-indices $\mi{e}_i$\\
\code{NUM\_alpha}, \code{HOD\_D0}, \code{HOD\_KMAX} & the carried count and the $D_0$, $K$ of construction\\
\code{hod\_tables\_ready} & whether the shared tables are built\\
\code{tanh\_derivs\_hod}, \code{tanh\_derivs\_ts} & $\tanh^{(q)}$ for $q=0..K{+}1$ by recursion\\
\code{write\_alpha\_order(file)} & write the multi-index ordering to a file\\
\code{bell\_*}, \code{pair\_*}, \code{fq\_*} & the Bell, adjoint and forward coefficient tables (internal structure, but public)\\
\bottomrule
\end{longtable}

\subsection{\code{net\_module}: networks and evaluation}

\begin{longtable}{@{}p{0.36\textwidth}p{0.58\textwidth}@{}}
\toprule
Element & Role \\
\midrule
\endhead
\code{net\_t} & the network type (layer count, widths, weights, a reference to its table set)\\
\code{net\_init(nt,L,ndim,ts)} & construction; \code{ts} decides the carried multi-index set\\
\code{net\_load}, \code{net\_save} & read/write in the \code{nn\_weight.dat} format\\
\code{net\_free} & release; derived types are not auto-released through \code{intent(OUT)}, so always call it\\
\code{work\_t}, \code{work\_init}, \code{work\_free} & the evaluation work space, one per thread\\
\code{net\_eval\_hod(nt,wk,x,t)} & one point's value and all carried derivatives (the inference path)\\
\bottomrule
\end{longtable}

\subsection{\code{train\_module}: training}

\begin{longtable}{@{}p{0.36\textwidth}p{0.58\textwidth}@{}}
\toprule
Element & Role \\
\midrule
\endhead
\code{twork\_t}, \code{twork\_init}, \code{twork\_free} & the adjoint work space (holds the forward pass's intermediates)\\
\code{grad\_t}, \code{grad\_init}, \code{grad\_free} & the gradient accumulator; optimizer moments live here too\\
\code{grad\_zero}, \code{grad\_add} & zeroing and addition; \code{grad\_add} is the thread-parallel reduction\\
\code{net\_forward\_point(nt,tw,x,t)} & one point's forward pass; returns all carried derivatives and keeps the adjoint intermediates\\
\code{net\_backward\_point(nt,tw,seed,g)} & the adjoint of a seed, accumulated into \code{g}\\
\code{net\_grad\_point(nt,tw,x,seed,g)} & the two above in one call\\
\code{bwork\_t}, \code{bwork\_init}, \code{bwork\_free} & the batch work space (carries a batch axis)\\
\code{net\_grad\_batch(nt,bw,X,Y,c,g)} & the batch gradient of a value fit; three matrix products per layer\\
\code{opt\_sgd\_step}, \code{opt\_adam\_step} & one step applied to a caller-owned array\\
\bottomrule
\end{longtable}

\subsection{\code{kalman\_module}: extended-Kalman-filter training}

A training route separate from gradient descent, used directly from
the library or a host code (the application's Kalman option in
\code{GD\_method} is the trainer-side counterpart). It is the method
in standard use for high-dimensional neural-network potentials of the
n2p2 family.

One update presents one scalar observable (an energy, or one force
component the host has chained into descriptor space) and applies the
rank-one form; no matrix is inverted:
\begin{align*}
\xi &= y_{\rm obs} - \hat y, &
\mi{P}\mi{j} &= \mi{P}\,\mi{j}, &
\mi{K} &= \frac{\mi{P}\mi{j}}{\lambda + \mi{j}^{\!\top}\mi{P}\mi{j}},\\
\mi{w} &\leftarrow \mi{w} + \mi{K}\xi, &
\mi{P} &\leftarrow \frac{\mi{P}-\mi{K}(\mi{P}\mi{j})^{\!\top}}{\lambda}, &
\lambda &\leftarrow \lambda_0\lambda + 1 - \lambda_0 .
\end{align*}
The Jacobian row $\mi{j}=\partial\hat y/\partial\mi{w}$ comes from
\code{net\_grad\_point} with a unit seed on the observed slot, so
\textbf{any carried derivative can be the observable}, Hessian
components included.

The covariance $\mi{P}$ is quadratic in the number of weights. That is
the method's intended regime: a small network updated very few times
per pattern (the usual trade against SGD/Adam).

\begin{longtable}{@{}p{0.36\textwidth}p{0.58\textwidth}@{}}
\toprule
Element & Role \\
\midrule
\endhead
\code{kalman\_t} & the filter type ($\mi{P}$, forgetting factor $\lambda$, $\lambda_0$)\\
\code{kf\_num\_weights(nt)} & the state dimension, for sizing $\mi{P}$\\
\code{kf\_init(kf,nt,P0,lam,lam0)} & construction; $\mi{P}_0$ is the diagonal initial value\\
\code{kf\_update(kf,nt,tw,g,x,islot,y)} & update on the observation of slot \code{islot} at point \code{x} with value \code{y}\\
\code{kf\_update\_obs(kf,nt,j,xi)} & the form taking the Jacobian row and residual directly, for callers that assemble their own chain rule\\
\code{kf\_free} & release\\
\bottomrule
\end{longtable}

As verification, \code{libverify\_example} confirms that Kalman
training on a linear network reproduces the ridge-regression
normal-equation solution to machine precision (measured
$5\times10^{-13}$).

\subsection{\code{committee\_module}: uncertainty by committee}

Returns the spread of several networks' predictions at the same input
as an uncertainty indicator. The intended use trains the members
separately (for instance with different \code{Rand\_seed}) and bundles
them at evaluation.

\begin{longtable}{@{}p{0.36\textwidth}p{0.58\textwidth}@{}}
\toprule
Element & Role \\
\midrule
\endhead
\code{committee\_t} & the committee type (a set of member references)\\
\code{comm\_init}, \code{comm\_free} & construction and release\\
\code{comm\_eval} & mean and spread at one point\\
\bottomrule
\end{longtable}

\subsection{\code{blas\_wrap\_module}}

A single entry, \code{bgemm}: the $\alpha=1,\beta=0$ matrix product
behind an explicit interface. The library's work spaces are pointer
components, whose elements cannot be sequence-associated to the
assumed-size dummies of the implicit BLAS interface, so callers hand
whole arrays to this wrapper. Built without \code{USE\_BLAS} it falls
back to reference loops.

\subsection{The C bindings}
\label{sec:capi}

\code{lib/c\_api\_module.f90} and \code{lib/dnnf90.h} provide 22
functions. They avoid the \code{.mod} incompatibility, so embeddings
that cross toolchains use these.

\begin{longtable}{@{}p{0.42\textwidth}p{0.52\textwidth}@{}}
\toprule
Function group & Role \\
\midrule
\endhead
\code{dnnf90\_tables\_init\_dense}, \code{dnnf90\_tables\_init\_closure} & shared tables (dense / closure)\\
\code{dnnf90\_tabset\_dense}, \code{dnnf90\_tabset\_closure}, \code{dnnf90\_tabset\_nderiv} & table sets and carried-count query\\
\code{dnnf90\_net\_create}, \code{dnnf90\_net\_create\_ts}, \code{dnnf90\_net\_set\_layer}, \code{dnnf90\_net\_load}, \code{dnnf90\_net\_free} & network construction, layer setup, loading, release\\
\code{dnnf90\_work\_create}, \code{dnnf90\_work\_free}, \code{dnnf90\_twork\_create}, \code{dnnf90\_twork\_free} & work spaces (evaluation / adjoint)\\
\code{dnnf90\_grad\_create}, \code{dnnf90\_grad\_zero}, \code{dnnf90\_grad\_free} & gradient accumulators\\
\code{dnnf90\_eval} & one-point evaluation\\
\code{dnnf90\_grad\_point} & one point's weight gradient from a seed\\
\code{dnnf90\_adam\_step} & one step\\
\code{dnnf90\_nderiv}, \code{dnnf90\_alpha} & carried-count and multi-index queries\\
\bottomrule
\end{longtable}

C wrappers for the committee and the Kalman filter are not provided
(Section~\ref{sec:limits}).

%=====================================================================
\section{Input specification (\code{input\_nn.dat})}

\subsection{Format conventions}

\begin{itemize}[leftmargin=1.4em]
\item One key per line, \code{KEY value ... /}; everything after
      \code{/} is ignored as a remark.
\item A leading \code{\#} or \code{!} makes a comment line.
\item Key names are case-insensitive, and so are the named values
      (\code{Task}, \code{Activation}, \code{Init\_w},
      \code{GD\_method}, the \code{Loss\_term} form names, the
      residual term names, the NGD schedule names).
\item If a key appears more than once, \textbf{the later line wins}.
\item Key order is free. The only exceptions are the keys that need
      \code{Nlayer} first (\code{Loss\_term}, \code{Residual},
      \code{Time\_axis}); appearing before \code{Nlayer} they stop
      the run.
\item A word not in the parser's branches is skipped with the warning
      \code{input WARNING: unrecognized keyword ignored:}. A retired
      key is \emph{not} skipped: it stops the run and states its
      migration path.
\end{itemize}

\subsection{Required keys and defaults}

Only five keys are required: \code{Nlayer}, \code{Init\_w},
\code{Loss\_term}, \code{GD\_method}, \code{GD\_param}. High-order
fitting or collocation adds \code{Hod\_K} and \code{Residual}.
Everything else has a default. The smallest input is:

\begin{verbatim}
Init_w     Glorot_norm
Nlayer     3 /
1 /
8 /
1 /
Loss_term  DATA  200  train.dat  1.d0 /
Epoch      200 /
GD_method  Adam
GD_param   0.01d0 0.9d0 0.999d0 1.d-8 0.d0 /
\end{verbatim}

\subsection{Network and training}

\begin{longtable}{@{}p{0.27\textwidth}p{0.13\textwidth}p{0.54\textwidth}@{}}
\toprule
Key & Default & Meaning \\
\midrule
\endhead
\code{Nlayer L /} & required & layer count; the $L$ lines that follow give one width each. Layer 1 is the input dimension $D_0$; the last layer must be 1\\
\code{Activation TANH /} & \code{TANH} & hidden activation. Anything but \code{TANH} stops the run, since the derivative tables are built from the Bell polynomials of $\tanh$ (see Section~\ref{sec:activation} for the other tabulated activations)\\
\code{Init\_w NAME} & required & \code{Glorot\_unif}, \code{Glorot\_norm}, \code{He\_unif}, \code{He\_norm}, \code{LeCun\_unif}, \code{LeCun\_norm}, \code{Random\_unif}, \code{Random\_norm}. The first six set the variance from the fan-in (or fan-in and fan-out). \code{Random\_*} is a fixed width-independent range ($\pm 0.05$, or $\sigma=0.05$), under which deep networks lose both signal and gradient\\
\code{Rand\_seed n /} & 0 (clock) & random seed; always set it for reproducibility\\
\code{Restart n /} & 0 & 1 resumes from the weights in \code{nn\_weight.dat}; the optimizer state is also restored when \code{gd\_*.dat} exists (Section~\ref{sec:restart})\\
\code{Task TRAIN|PREDICT|COMMITTEE /} & \code{TRAIN} & \code{PREDICT} reads \code{nn\_weight.dat}, writes the output files and exits. \code{COMMITTEE} evaluates a set of trained networks and exits (Section~\ref{sec:committee}). Neither produces history, checkpoints or optimizer logs\\
\code{Epoch n /} & 1000 & maximum epoch count\\
\code{Num\_batch n /} & $\min(10,N_{\rm train})$ & minibatch size; outside $[1,N_{\rm train}]$ the run stops\\
\code{Shuffle n /} & 0 & 1 redraws the training points each epoch\\
\code{Num\_validation n /} & 0 & points held out for validation; outside $[0,N_{\rm input}-1]$ the run stops. At 0 the validation cost is undefined and early stopping reads the training cost\\
\code{Validation\_cyc n /} & 1 & validation cadence (epochs); the history holds one line per validation\\
\code{P\_max n /} & effectively infinite & the early-stopping patience: the run stops after this many validations without improvement. By default there is no early stop\\
\code{Conv x /} & 0 (off) & stop when the training cost falls below this\\
\code{Checkpoint\_cyc n /} & 1000 & cadence of \code{checkpoint\_ep*.dat}\\
\code{Output\_deriv n /} & 0 & 1 writes the model's first derivatives after training; in a value-fit-only run the dense $K=1$ table is built automatically\\
\code{Hod\_check n /} & 0 & 1 runs the in-run self-check before training (Section~\ref{sec:verify})\\
\code{Average\_cyc\_mpi n /} & 1 & weight-averaging cadence under MPI\\
\code{Shuffle\_cyc\_mpi n /} & $10^6$ & training-point redistribution cadence under MPI\\
\bottomrule
\end{longtable}

\subsection{Choosing the activation}
\label{sec:activation}

The adjoint reaches order $K+1$, so the activation must be
$C^{K+1}$. Among functions that satisfy this, \textbf{how fast the
high-order derivatives grow} differs greatly, and this directly
controls the conditioning of the carried slots. Measuring
$\max_a|\sigma^{(q)}(a)|$ over $a\in[-6,6]$:

\begin{center}
\begin{tabular}{@{}lrrrrl@{}}
\toprule
Activation & $q=2$ & $q=4$ & $q=6$ & $q=8$ & parity \\
\midrule
\code{TANH}    & 0.77 & 4.0  & 50   & 1220  & odd \\
\code{ERF}     & 0.96 & 4.4  & 36   & 454   & odd \\
\code{SIN}     & 1.0  & 1.0  & 1.0  & 1.0   & odd \\
\code{BESSEL}  & 0.50 & 0.375& 0.31 & 0.273 & even \\
\code{BESSEL1} & 0.41 & 0.334& 0.288& 0.257 & odd \\
\bottomrule
\end{tabular}
\end{center}

\code{SIN} has $\sigma^{(q)}(a)=\sin(a+q\pi/2)$, bounded by one at
every order. \code{BESSEL} ($J_0$) and \code{BESSEL1} ($J_1$) have
derivatives given by the binomial combination
\[
J_\nu^{(q)}(a)=\frac{(-1)^q}{2^q}\sum_{k=0}^{q}(-1)^k\binom{q}{k}J_{\nu+q-2k}(a)
\]
and are likewise bounded, since $|J_n|\le1$ at every order. $J_0$ is
even, so the layer loses the sign of its preactivation; $J_1$ is odd
with the same boundedness.

Only the \code{BESSEL} family uses externally derived code: the four
routines \code{jyna} etc.\ of Zhang and Jin live in
\code{lib/bessel\_module.f90} with their copyright notice preserved
(see \code{LICENSE}). The original routines treat negative arguments
as $x=0$, so a wrapping layer applies the reflection
$J_n(-x)=(-1)^nJ_n(x)$ --- preactivations are of course signed, and
omitting this would make half the layer return the constant $J_0=1$.

For \code{ERF}, only the value uses the Fortran 2008 intrinsic
\code{ERF}, and that single file is excluded from
\code{make f95check}. The derivatives are the Hermite functions
$\mathrm{erf}^{(q)}(a)=\frac{2}{\sqrt\pi}(-1)^{q-1}H_{q-1}(a)e^{-a^2}$,
carried by the three-term recursion $\psi_{n+1}=2a\psi_n-2n\psi_{n-1}$
(with $\psi_n=H_ne^{-a^2}$), so no bare polynomial is formed and
nothing overflows.

The oscillatory activations (\code{SIN}, \code{BESSEL},
\code{BESSEL1}) can use \code{Init\_w SIREN}: an initialization that
places the preactivations inside one period, with
\code{Init\_w\_omega} (default 30) as the frequency. The optimum
depends on the input range; on the shipped example with inputs in
$[-1,1]$, 12 works and 30 does not.

\subsection{Evaluating a high-order fit}
\label{sec:hodmetric}

A single coefficient of determination pooled over all carried slots
is useless at large $K$. On the shipped $K=7$ example (330 slots) the
mean square of the targets is nearly flat across orders
($0.04$--$0.07$) while the predictions blow up at high order, so one
number hides the structure.

\code{bench/post/hod\_accuracy.py} reports two dimensionless
quantities per order,
\[
\mathrm{rel}(p)=\frac{\|T_{\rm pred}-y\|_p}{\|y\|_p},\qquad
\mathrm{amp}(p)=\frac{\|T_{\rm pred}\|_p}{\|y\|_p},
\]
where $\mathrm{rel}\ge1$ means the order carries no information and
$\mathrm{amp}\gg1$ means it diverges. Rankings use the relative error
in the loss norm, $\|T-y\|_\lambda/\|y\|_\lambda$.

Measured this way, the activations differ structurally at $K=7$:
\code{TANH} reaches $\mathrm{amp}=14.7$ at order seven --- the
prediction swings fifteen times the target --- against 4.7 for
\code{SIN}, 5.2 for \code{BESSEL}, 3.9 for \code{BESSEL1}.

\paragraph{What $\lambda_p$ does.}
The shipped example uses $\lambda_p=1/(p!)^2$, so order seven carries
weight $4\times10^{-8}$: the high orders are barely optimized at all.
Re-measuring at $\lambda_p=1$ (equal weight across orders) changes the
picture entirely.

\begin{center}
\begin{tabular}{@{}lrrr@{}}
\toprule
 & \multicolumn{2}{c}{$\mathrm{rel}(7)$} & rel.\ error, loss norm \\
Activation & $\lambda_p=1/(p!)^2$ & $\lambda_p=1$ & $\lambda_p=1$ \\
\midrule
\code{TANH}   & 14.57 & 1.05 & 1.010 \\
\code{SIN}    &  4.33 & 0.45 & 0.377 \\
\code{BESSEL} &  5.04 & 0.31 & 0.221 \\
\bottomrule
\end{tabular}
\end{center}

Equal weighting tames the order-seven divergence for every
activation, but \code{TANH} then sits at $\mathrm{rel}\approx1$
\textbf{at every order}, with $\mathrm{amp}$ collapsed to
$0.05$--$0.47$: asked to actually deliver high orders, the $\tanh$
network degenerates to the mean everywhere. \code{BESSEL} holds
$\mathrm{rel}=0.12$--$0.31$ and $\mathrm{amp}\approx0.9$ at every
order.

``The high orders do not fit'' is therefore not an activation problem
alone but also a loss-design problem --- $\lambda_p$ may be throwing
the high orders away. In settings that genuinely demand them, the
boundedness of the activation's high-order derivatives is decisive.

\subsection{Loss terms and their data files}

\subsubsection{The form of the loss}

The loss is a weighted sum of terms. One \code{Loss\_term} line
declares one term; terms are numbered in order of appearance (the
number that appears in the output file names).
\[
L=\sum_{t} w_t\, L_t,\qquad
L_t=\frac{1}{2}\sum_{n\in t}\begin{cases}
\big(N(\mi{x}_n)-y_n\big)^2 & \text{DATA}\\[2pt]
\displaystyle\sum_{\mi\alpha}\lambda_{|\mi\alpha|}\big(T^{\mi\alpha}(\mi{x}_n)-y_{n,\mi\alpha}\big)^2 & \text{HOD\_DATA}\\[2pt]
R(\mi{x}_n)^2 & \text{COLLOCATION}
\end{cases}
\]
$w_t$ is the fourth argument of the \code{Loss\_term} line.
$\lambda_{|\mi\alpha|}$ is the \textbf{per-order} weight, and
$\lambda_0$ multiplies the value itself ($\mi\alpha=\mi 0$): a run
with a single \code{HOD\_DATA} term already contains a value fit
whenever $\lambda_0>0$. Conversely, $\lambda_0=0$ leaves the value
unconstrained --- a function whose derivatives fit but whose constant
is free: on \code{zwork/hod\_4d\_k3}, switching $\lambda_0$ from 1 to
0 drops the value's coefficient of determination from $0.985$ to
$-17.1$.

\subsubsection{What each form requires}

This table narrows down what a given problem needs.

\begin{center}
\begin{tabular}{@{}llll@{}}
\toprule
Form & Required keys & Data columns & Use \\
\midrule
\code{DATA} & none & $D_0+1$ & supervised values, boundary conditions \\
\code{HOD\_DATA} & \code{Hod\_K} & $D_0+N_\alpha$ & supervised derivatives \\
\code{COLLOCATION} & \code{Residual} (\code{Hod\_K} optional) & $D_0$ + extra columns & points where the equation is imposed \\
\bottomrule
\end{tabular}
\end{center}

$D_0$ is the number of input variables (the first width after
\code{Nlayer}); $N_\alpha$ the number of carried multi-indices
(Section~\ref{sec:nalpha}). Data files are read list-directed, so
\textbf{records are numbers only} --- comment lines are not skipped.
A record may wrap across lines (tokens are counted). Lines beyond the
\code{n} of \code{Loss\_term} are not read.

\subsubsection{Worked examples}

\paragraph{A value fit (\code{DATA}).}
$D_0=1$, $y=x^2+1$ at 200 points; one record is the two numbers $x$
and $y$.

\begin{verbatim}
--- input_nn.dat ---            --- train.dat (200 lines) ---
Nlayer     4 /                  -0.99004975124378  1.98019850993787
1 /   <- D0 = 1                 -0.98009950248756  1.96059503477637
12 /                            ...
12 /
1 /
Loss_term  DATA  200  train.dat  1.d0 /
\end{verbatim}

\paragraph{A derivative fit (\code{HOD\_DATA}).}
$D_0=4$, $K=3$, dense set, so $N_\alpha=35$. One record is
$4+35=39$ numbers: $\mi{x}$ first, then the 35 targets in multi-index
order. Note that the $K+1$ lines directly after \code{Hod\_K} are
$\lambda_0,\dots,\lambda_K$ (they are continuation lines belonging to
\code{Hod\_K}; no other key may be interposed).

\begin{verbatim}
--- input_nn.dat ---
Nlayer     4 /
4 /   <- D0 = 4
10 /
10 /
1 /
Hod_K      3 /      <- K = 3 ; the next K+1 = 4 lines are the lambdas
1.d0                <- lambda_0 : multiplies the value
1.d0                <- lambda_1 : multiplies the first derivatives
1.d0                <- lambda_2
1.d0                <- lambda_3
Hod_dense  1 /      <- the dense set |alpha| <= K
Loss_term  HOD_DATA  400  train.dat  1.d0 /

--- train.dat (400 records, 39 numbers each) ---
 0.989957667989E+00 -0.774791878492E-01 -0.168822326093E+00  0.269057713255E+00  ...
 |<---------------- x(1:4) ---------------->|<-- y_alpha(1:35) ...
\end{verbatim}

\paragraph{Points where the equation is imposed (\code{COLLOCATION}).}
$D_0=2$. The residual contains \code{SRC} and
\code{Exact\_solution 1} is set, so one record is $\mi{x}$ (2
numbers), the source $f$, and the exact solution $u^*$, \textbf{in
that order} --- 4 numbers. The exact solution is for reporting only;
it does not enter the loss.

\begin{verbatim}
--- input_nn.dat ---
Loss_term  DATA         160  data.dat     1.d0  /  <- boundary conditions
Loss_term  COLLOCATION  400  colloc.dat   1.d-3 /  <- interior points
Residual 3 /
LIN   1.d0    2 0 /
LIN   1.d0    0 2 /
SRC  -1.d0        /   <- declares that colloc.dat has an f column
Exact_solution 1 /    <- declares that colloc.dat has a u* column

--- colloc.dat (400 records, 4 numbers each) ---
  6.91978764042226e-01  9.36279743968240e-01 -3.23261059769365e+00  1.63765966006543e-01
  |<------- x(1:2) ------->|<---- f(x) ---->|<---- u*(x) ---->|

--- data.dat (160 records, 3 numbers each: DATA, so x and y) ---
  1.25000000000000e-02  0.00000000000000e+00  0.00000000000000e+00
  |<------- x(1:2) ------->|<--- y = 0 --->|
\end{verbatim}

The extra columns exist \textbf{only when declared}. Without
\code{SRC} there is no $f$ column; with \code{Exact\_solution 0}
(the default) no $u^*$ column. The smallest collocation file is the
$D_0$ columns of $\mi{x}$ alone.

\subsubsection{The carried multi-index set and $N_\alpha$}
\label{sec:nalpha}

Writing a \code{HOD\_DATA} file requires knowing $N_\alpha$ and the
\textbf{column order} first. The set is decided as follows.

\begin{center}
\begin{tabular}{@{}lll@{}}
\toprule
Condition & Carried set & $N_\alpha$ \\
\midrule
\code{Hod\_dense 1} & all $|\mi\alpha|\le K$ & $\binom{D_0+K}{K}$ \\
\code{Hod\_alpha\_file} given & the downward closure of the seed set & decided at run time \\
\code{COLLOCATION} only & the closure of the residual's multi-indices & decided at run time \\
\bottomrule
\end{tabular}
\end{center}

For the dense set, $N_\alpha$ is computable in advance ($D_0=4,K=3$
gives $\binom{7}{3}=35$). The order is \textbf{ascending in
$|\mi\alpha|$}, lexicographic within an order. Every run writes
\code{hod\_alpha\_order.dat}, which is the definition of the column
order.

\begin{verbatim}
--- hod_alpha_order.dat  (D0 = 4, K = 3) ---
# ia  |alpha|  alpha(1:D0)
1  0   0 0 0 0      <- the value
2  1   1 0 0 0      <- d/dx1
3  1   0 1 0 0      <- d/dx2
4  1   0 0 1 0
5  1   0 0 0 1
6  2   2 0 0 0      <- d^2/dx1^2
...
\end{verbatim}

\paragraph{The theory of closures.}
A multi-index set $\mathcal{A}$ is \textbf{downward closed} when
$\mi\alpha\in\mathcal{A}$ and $\mi\beta\le\mi\alpha$ (componentwise)
imply $\mi\beta\in\mathcal{A}$. The carried set must be closed
because of the recursion's structure: the Bell part of
Section~\ref{sec:fwd} builds $T^{\mi\alpha}$ from $S^{\mi\beta}$ and
$B_{\mi\alpha-\mi\beta,q-1}$ with $\mi\beta\le\mi\alpha$, and the
adjoint reads $h_{\mi\alpha-\mi\beta}$; an index whose sub-indices
are absent cannot be computed at all.

The closure of a seed set $\mathcal{S}$ is
\[
\mathrm{cl}(\mathcal{S})=\bigcup_{\mi\alpha\in\mathcal{S}}
\mathrm{box}(\mi\alpha),\qquad
\mathrm{box}(\mi\alpha)=\{\mi\beta:\mi\beta\le\mi\alpha\},\qquad
|\mathrm{box}(\mi\alpha)|=\prod_{i=1}^{D_0}(\alpha_i+1),
\]
a union of boxes, so
$|\mathrm{cl}(\mathcal{S})|\le\sum_{\mi\alpha}\prod_i(\alpha_i+1)$
with equality for disjoint boxes; the exact count follows by
inclusion--exclusion.

Take the shipped \code{zwork/hod\_4d\_k7\_active} ($D_0=4$). The
seeds are $(7,0,0,0)$, $(0,0,0,1)$, $(2,2,1,0)$, with box sizes $8$,
$2$, $3\cdot3\cdot2=18$. The intersection
$\mathrm{box}(7,0,0,0)\cap\mathrm{box}(2,2,1,0)=\{(k,0,0,0):k\le2\}$
has $3$ elements; the other pairwise and the triple intersections are
all $\{\mi 0\}$ with $1$. Hence
\[
|\mathrm{cl}(\mathcal{S})|=8+2+18-3-1-1+1=24 .
\]
The run reports exactly $N_\alpha=24$ --- under 7\% of the dense
$\binom{4+7}{7}=330$. Since the cost is $O(L\,n^2 N_\alpha)$, that
ratio is exactly the gain in speed and memory a closure buys.

In a \code{COLLOCATION}-only run, no seed file is needed:
\textbf{the residual's multi-indices are the seeds automatically}.
The \code{bench/kdv} residual contains $u_t$, $u\,u_x$, $u_{xxx}$,
i.e.\ $(0,1)$, $(1,0)$, $(3,0)$, whose closure is the five indices
$\{(0,0),(1,0),(2,0),(3,0),(0,1)\}$ --- half the dense
$\binom{5}{3}=10$, as the run also reports.

\paragraph{Ordering.}
The column order is ascending order of $|\mi\alpha|$, and within an
order it is the generation order of non-decreasing axis combinations
(with repetition): choosing axis $i$ once adds one to $\alpha_i$, so
at order 2 with $D_0=4$ the tuples
$(1,1),(1,2),(1,3),(1,4),(2,2),\dots$ produce
$\mi\alpha=(2,0,0,0),(1,1,0,0),(1,0,1,0),(1,0,0,1),(0,2,0,0),\dots$
A closure \textbf{keeps that generation order and drops the absent
indices}, so a closure preserves the dense set's relative order.
Indexing is a binary search on the lexicographic order
(\code{find\_alpha}), which this ordering underwrites.

\paragraph{Getting the column order without a run.}
\code{tools/alpha\_order.py} reproduces the same enumeration, byte
for byte against \code{hod\_alpha\_order.dat}, so data files can be
prepared before any run.

\begin{verbatim}
# the dense set
tools/alpha_order.py --d0 4 --k 3

# a closure (a seed file in the Hod_alpha_file format)
tools/alpha_order.py --d0 4 --seeds zwork/hod_4d_k7_active/alpha_seeds.dat

# just the count
tools/alpha_order.py --d0 10 --k 3 --count      ->  286
\end{verbatim}

The three shipped variants (dense $D_0{=}4,K{=}3$, a seed closure, a
residual-derived closure) are verified byte-identical against the
run's \code{hod\_alpha\_order.dat}. Running once to obtain the order
therefore still works, but is not necessary.

A record with the wrong column count is caught at read time; the run
stops stating the measured and the required counts:

\begin{verbatim}
 read_data: record 1 of train.dat
   accumulates 59 values where ndim(1)+NUM_alpha = 39 are required:
   a line is missing a value
   (x(1:ndim(1)) then y_alpha for every carried multi-index)
\end{verbatim}

Because tokens are counted, a short record \textbf{eats into the next
line} before the mismatch is noticed; the ``measured'' number is the
accumulated count after the boundary slipped, not the token count of
one line.

\subsubsection{Key list}

\begin{longtable}{@{}p{0.32\textwidth}p{0.62\textwidth}@{}}
\toprule
Key & Meaning \\
\midrule
\endhead
\code{Loss\_term FORM n file w /} & declare one loss term. \code{FORM} is \code{DATA}, \code{HOD\_DATA} or \code{COLLOCATION}; \code{n} the record count to read, \code{file} the data file, \code{w} the term weight ($w\ge 0$; negative stops). At most 100 terms\\
\code{Hod\_K K /} & the maximum carried order $K$. \textbf{The next $K+1$ lines} hold $\lambda_0,\dots,\lambda_K$, one per line. Required for \code{HOD\_DATA}; optional for \code{COLLOCATION}-only runs, where the residual's highest order becomes $K$. An explicit $K$ smaller than the residual's order stops the run\\
\code{Hod\_dense n /} & 1 forces the dense set $|\mi\alpha|\le K$. At 0 (default) the closure is built from the \code{Hod\_alpha\_file} seeds and the residual's multi-indices; however, a \code{HOD\_DATA} term without \code{Hod\_alpha\_file} implies the dense set, since the data columns correspond to all multi-indices\\
\code{Hod\_alpha\_file f /} & the seed multi-index file: first line the count, then one $\mi\alpha(1{:}D_0)$ per line\\
\bottomrule
\end{longtable}

\subsection{Residuals (collocation)}

\begin{longtable}{@{}p{0.32\textwidth}p{0.62\textwidth}@{}}
\toprule
Key & Meaning \\
\midrule
\endhead
\code{Residual m /} & define the residual operator in the next $m$ lines, each one of the four kinds below. At most 64 terms after expansion\\
 & \quad \code{LIN c a1..aD /} \quad $c\,T^{\mi\alpha}$ (a linear term)\\
 & \quad \code{UUX c a1..aD /} \quad $c\,u\,T^{\mi\alpha}$ ($c\,u^2$ when $\mi\alpha=\mi0$)\\
 & \quad \code{DXLAP c k [ix] /} \quad $c\,\partial_{ix}\nabla^{2k}u$. The Laplacian runs over every non-time axis and expands multinomially into several linear terms. \code{ix} defaults to axis 1; a time-axis \code{ix} stops the run. At most 16 \code{DXLAP} lines\\
 & \quad \code{SRC c /} \quad $c\,f(\mi{x})$ (the source, read per collocation point; at most once)\\
\code{Time\_axis n /} & the index of the time axis; default the last axis, and every shipped case follows that convention. \code{DXLAP} expands after the whole input is read, so this key may appear anywhere\\
\code{Exact\_solution n /} & 1 reads the collocation file's last column as the exact solution and adds the error to the output; it never enters the loss\\
\bottomrule
\end{longtable}

The residual is evaluated pointwise as
$R=\sum_k c_k\,[u]\,T^{\mi\alpha_k}+c_{\rm src}f$, and the seed is
$\partial L/\partial T^{\mi\alpha}= w\,R\,\partial R/\partial
T^{\mi\alpha}$; the source is additive and does not appear in
$\partial R/\partial T^{\mi\alpha}$. Residual evaluation and seed
construction are argument-only functions with no shared state.

A collocation run writes the parsed residual back to the log head in
human-readable form. The operator is the one quantity in the input
with nothing to cross-check against --- a mistyped multi-index just
solves a different equation correctly --- so this echo is the only
defence:

\begin{verbatim}
### Residual as parsed (imposed as R = 0 at every collocation point)
###   axes are written x1..x2, with axis 2 shown as t (Time_axis)
###   term      coefficient  expression                    slot  multi-index
###    1      1.000000E+00  u_t                               3    0  1
###    2      3.000000E+00  u * u_x1                          2    1  0
###    3      1.000000E+00  u_x1x1x1                          5    3  0
\end{verbatim}

The last column is the slot of the carried multi-index, checkable
against \code{hod\_alpha\_order.dat}.

\subsection{Coupled systems (the \code{System} block)}
\label{sec:system}

The scalar \code{Residual} block writes one residual of one field.
\textbf{Systems of residuals over several field components} ---
incompressible Navier--Stokes, or a charge density coupled to a
potential --- are written with the \code{System} block. The output
width $m$ becomes the number of field components (the scalar form's
width-1 constraint is lifted when \code{System} is present), and
training carries every derivative of every component to order $K$.

Each term is a product of at most four derivative factors,
\[
  c\;\partial^{\mi\gamma}u_{k}\;\partial^{\mi\beta}u_{i}\;\partial^{\mi\alpha}u_{j},
\]
and the four term kinds are its cases: \code{TRM} one factor,
\code{XUX} two with the first undifferentiated ($\mi\beta=\mi0$),
\code{DXD} two differentiated factors, \code{TRP} three. With one
component and \code{ic = jc = 1} they reduce to \code{LIN}/\code{UUX},
so existing inputs keep their meaning.

\begin{longtable}{@{}p{0.46\textwidth}p{0.48\textwidth}@{}}
\toprule
Key & Meaning \\
\midrule
\endhead
\code{System nres nterm /} & define \code{nres} residuals ($\le 16$) in the next \code{nterm} lines ($\le 256$). \code{ir} is the residual a term belongs to; \code{jc}, \code{ic}, \code{kc} are component indices\\
\code{TRM ir jc a1..aD c /} & $c\,\partial^{\mi\alpha}u_{jc}$ (linear)\\
\code{XUX ir ic jc a1..aD c /} & $c\,u_{ic}\,\partial^{\mi\alpha}u_{jc}$: advection $u\,u_x+v\,u_y$, charge transport by a flow, the body force $\rho E$ --- the terms coupled systems need and the scalar form cannot write\\
\code{DXD ir ic b1..bD jc a1..aD c /} & $c\,\partial^{\mi\beta}u_{ic}\,\partial^{\mi\alpha}u_{jc}$ (a product of derivatives). Flux divergences need it: in $\nabla\!\cdot\!(\rho\nabla\phi)=\nabla\rho\!\cdot\!\nabla\phi+\rho\,\nabla^2\phi$ the first term is a product of first derivatives of different components\\
\code{TRP ir kc g1..gD ic b1..bD jc a1..aD c /} & $c\,\partial^{\mi\gamma}u_{kc}\,\partial^{\mi\beta}u_{ic}\,\partial^{\mi\alpha}u_{jc}$ (three factors), needed by the compressible momentum flux $\rho u u_x$ or field-dependent mobilities\\
\code{QAD ir lc d1..dD kc g1..gD ic b1..bD jc a1..aD c /} & $c\,\partial^{\mi\delta}u_{lc}\,\partial^{\mi\gamma}u_{kc}\,\partial^{\mi\beta}u_{ic}\,\partial^{\mi\alpha}u_{jc}$ (four factors), the quartic nonlinearity of the higher dispersive hierarchies: the seventh-order Lax equation carries $140\,u^3u_x$\\
\code{Sys\_src /} & declares that the collocation file carries \textbf{one source column per residual}, between the coordinates and the exact solutions. Omit for a homogeneous system, whose file is that much narrower; a mistaken setting stops with a record-length error --- the check working as intended\\
\code{Sys\_wcomp w1..wm /} & per-component weights of the supervised loss (default all 1). Mandatory for components of different magnitudes: an unweighted sum of squares is a statement about the largest component only. The custom is to measure $1/\lVert y_i\rVert^2$ from the data, as the shipped cases do\\
\code{Sys\_wres w1..wr /} & per-residual weights of the collocation loss (default all 1), for the same reason. Reads \code{nres} values, so it must follow \code{System}\\
\code{Sys\_rnoise r1..rr /} & per-residual observation noise for the extended Kalman filter. Exact, not approximate: scaling both the observation row and the innovation by $1/\sqrt{r}$ reproduces noise $r$ on the original quantity exactly. A residual with an outsized source produces outsized innovations, and a filter trusting all observations equally spends its covariance on that residual; raising its noise says ``this observation is less trustworthy''. Must follow \code{System}\\
\code{Sys\_balance cyc alpha /} & every \code{cyc} epochs, readjust \code{Sys\_wres} so that each residual contributes comparably to the gradient. Each residual's own gradient norm is sampled on up to 40 training points, and each weight moves a fraction $\alpha$ toward $g_{\max}/g_{ir}$ (the running average keeps it from chasing one batch's noise). The weights are printed at each readjustment\\
\bottomrule
\end{longtable}

The system data formats: \code{data.dat} is
$\mi{x}(1{:}D_0)\;y(1{:}m)$; \code{colloc.dat} is $\mi{x}(1{:}D_0)$,
then one source per residual if \code{Sys\_src}, then all components'
exact solutions if \code{Exact\_solution 1}. The reader counts
values, so a column slip stops with a record-length error rather than
propagating silently.

Residual evaluation (\code{calc\_sys\_residual}) and seed
construction (\code{set\_sys\_seed}) are argument-only functions,
like the scalar form. Each factor of a product contributes to the
seed by the product rule; the adjoint of that is checked against
finite differences for one to three factors by
\code{tools/example\_product\_adj.f90} ($9.5\times10^{-12}$), and for
the four-factor form by the in-run check of \code{zwork/pinn\_lax7}
($2.0\times10^{-10}$), and the
multi-output forward and adjoint by
\code{tools/example\_multiout*.f90} ($2.2\times10^{-10}$,
$9.1\times10^{-12}$).

The shipped examples are \code{zwork/pinn\_lax7} (the seventh-order Lax equation, whose quartic term needs the four-factor form; exact solitary wave, table read back to $1.1\times10^{-14}$), \code{zwork/pinn\_lax} (Lax's fifth-order
equation: a scalar case whose terms $30u^2u_x$ and $20u_xu_{xx}$ need
the product language; exact soliton, term table read back
independently to $8\times10^{-16}$),
\code{zwork/pinn\_kovasznay} (Kovasznay flow; exact solution, term
table verified against an independent implementation to $10^{-15}$),
\code{zwork/pinn\_cavity} (the Re$=100$ cavity, compared with the
Ghia benchmark profiles), and \code{zwork/pinn\_ehd} (a
five-component electrohydrodynamic system; manufactured solution).

\subsection{Uncertainty by committee}
\label{sec:committee}

\code{Task COMMITTEE} uses \code{committee\_module}
(Section~\ref{sec:libapi}) from the trainer. Members are
independently trained networks and cannot be produced in one run
(usually the same input is trained under different \code{Rand\_seed}
and each \code{nn\_weight.dat} saved under its own name). This task
reads them and writes, for every input point, \textbf{the mean and
sample standard deviation of every carried slot}, then exits. The
point is that the spread of the \emph{derivatives} is available, not
just of the value: with \code{Output\_deriv 1} the dense first-order
set is carried, so the spread of $\partial N/\partial x_i$ is
reported too --- exactly the quantity an on-the-fly scheme
thresholds when deciding whether to call the expensive
first-principles calculation.

\begin{longtable}{@{}p{0.30\textwidth}p{0.64\textwidth}@{}}
\toprule
Key & Meaning \\
\midrule
\endhead
\code{Committee m /} & the next $m$ lines name the members' weight files, one per line, $m\ge 2$. Members must share the multi-index set the current input defines; a slot-count mismatch stops the run\\
\bottomrule
\end{longtable}

The output is \code{output\_committee\_set<4 digits>.dat}: one line
is $\mi{x}(1{:}D_0)$, all slots' means, all slots' standard
deviations. At the end the spread over the whole input set is
summarized per slot (mean and maximum). The shipped example is
\code{zwork/tour\_committee} (\code{train\_members.sh} trains four
members; \code{committee.dat} evaluates them).

\subsection{Optimizers}

\begin{longtable}{@{}p{0.30\textwidth}p{0.64\textwidth}@{}}
\toprule
Key & Meaning \\
\midrule
\endhead
\code{GD\_method NAME} & the eleven gradient rules \code{SIMPLE}, \code{SIMPLE\_SCHEDULE}, \code{SIMPLE\_GMAXCLIP}, \code{MOMENTUM}, \code{NESTEROV}, \code{ADAGRAD}, \code{ADADELTA}, \code{RMSPROP}, \code{RMSPROP\_NESTEROV}, \code{ADAM}, \code{NATURAL\_GRAD}; the non-gradient \code{KALMAN}; and the quasi-Newton \code{LBFGS}\\
\code{GD\_param p1 ... p5 /} & the method's hyperparameters. An independent key, in either order with \code{GD\_method}. All five zero stops the run, since no method can take a step (only \code{LBFGS}, whose line search sets the step, is exempt and may omit \code{GD\_param}). The adaptive rules ($p_2$ for \code{ADAGRAD}, \code{ADADELTA}; $p_3$ for \code{RMSPROP}, \code{RMSPROP\_NESTEROV}; $p_4$ for \code{ADAM}) divide by an accumulator regularized by $\varepsilon$, so a non-positive slot stops the run (a weight with zero first-step gradient would divide by zero)\\
\bottomrule
\end{longtable}

$p_1$ is the learning rate for every rule. \code{SIMPLE\_SCHEDULE}
uses $\eta(t)=p_1/\sqrt{t+1}+p_2$; \code{MOMENTUM} and
\code{NESTEROV} take $p_2$ as the momentum; \code{ADAM} uses
$p_2,p_3,p_4$ as $\beta_1,\beta_2,\varepsilon$.

\code{KALMAN} selects the per-pattern extended Kalman filter of
\code{kalman\_module} (Section~\ref{sec:libapi}) from the trainer.
Rather than accumulating a gradient and stepping once, it presents
scalar observations one at a time and lets the rank-one update move
the weights itself. The observable per loss form: the value for
\code{DATA}; each carried derivative with
$\lambda_{|\mi\alpha|}\ne 0$ for \code{HOD\_DATA}; the residual
(target 0) for \code{COLLOCATION}. The Jacobian row
$\partial\hat y/\partial\mi{w}$ comes from the same seed path the
gradient rules use. \code{GD\_param} means: $p_1$ = the diagonal
initial value of $\mi{P}$; $p_2$ = the initial forgetting factor
$\lambda$; $p_3$ = the $\lambda_0$ of
$\lambda\leftarrow\lambda_0\lambda+1-\lambda_0$; the run stops unless
$p_1>0$ and $0<p_2,p_3\le 1$. $\mi{P}$ is quadratic in the weight
count, so this is for small networks; the run reports the size at
startup and refuses above 2\,GB.

\code{Kalman\_gate g /} is the innovation gate (a robust EKF): an
observation with $\xi^2 > g^2(\lambda + \mi{j}^{\top}\mi{P}\mi{j})$
has its row and innovation shrunk by the same factor (the same exact
mechanism as \code{Sys\_rnoise}) until equality --- every rank-one
step becomes bounded, at the price of absorbing surprising
observations more slowly. Default 0 (off).

The forgetting factor needs care. $\lambda<1$ re-inflates $\mi{P}$ by
$1/\lambda$ at every update, so under thousands of updates per epoch
the gains never decay and the filter can drift off a fitted state
along the linearization error of a nonlinear residual (observed on
coupled systems with source terms). With $p_2=p_3=1$ ($\lambda$
pinned at one) the gains decay as in recursive least squares and the
filter holds and refines a fitted state (a gentle descent at
$p_1=10^{-5}$). Forgetting buys initial convergence speed at the cost
of having no stationary behaviour: use forgetting for cold starts,
and $\lambda=1$ with small $p_1$ to refine a fitted state.
\code{Kalman\_iter n /} is the iterated EKF (IEKF): each observation
is relinearized $n$ times at the running iterate before the
covariance is updated once (Gauss--Newton on the single-observation
MAP problem). The measurement on the coupled EHD system is negative
and worth recording: neutral in the stable regime, and it
\textbf{amplifies} the divergence in the forgetting regime ($0.44$
against $0.023$ from the fitted start). Solving thousands of mutually
contradictory scalar observations each more exactly in sequence
strengthens the oscillation between them; the removed linearization
error had been acting as a buffer. The collective cure for this
failure mode is a simultaneous Gauss--Newton over all observations
--- that is, full-batch quasi-Newton.
\code{Kalman\_q q /} adds $q$ to the diagonal of $\mi{P}$ after each
update: process-noise injection, a stationary behaviour intermediate
between forgetting (inflating everything) and $\lambda=1$ (gains
vanishing). Default 0.
\code{Kalman\_mode DECOUPLED /} replaces the covariance with the
node-wise block diagonal (bias + fan-in as one block per neuron,
Puskorius--Feldkamp). The gain denominator
$\lambda+\mi{j}^{\top}\mi{P}\mi{j}$ stays global; only the
cross-neuron covariance is dropped. Memory and update cost fall from
$n_w^2$ to the sum of squared block sizes (measured 87$\times$ per
epoch on the shipped EHD net), the 2\,GB dense guard no longer
applies, and force-field widths become reachable. The price is the
convergence rate: on the KdV bench $1.2\times10^{-7}$ in 0.5\,s where
the dense filter reaches $7\times10^{-14}$ in about 6\,s; on Poisson
2D $3.2\times10^{-8}$ in 5.4\,s / 3000 epochs against
$5.5\times10^{-10}$ in 6.5\,s / 300 --- roughly ten times the epochs
per decade at a tenth the cost per epoch.

The shipped pair \code{zwork/pinn\_poisson2d} and
\code{zwork/pinn\_poisson2d\_kalman} differ only in the method (same
problem, data, network) and measure, against the exact solution:

\begin{center}
\begin{tabular}{@{}lrrrr@{}}
\toprule
Method & Epochs & Time & Training cost & RMSE \\
\midrule
\multicolumn{5}{@{}l}{\code{zwork/pinn\_poisson2d} (2-D Poisson: linear operator + source)}\\
Adam   & 20000 & 60 s & $1.1\times10^{-2}$ & $3.1\times10^{-3}$ \\
Kalman &   300 &  8 s & $4.8\times10^{-9}$ & $1.8\times10^{-5}$ \\
\midrule
\multicolumn{5}{@{}l}{\code{bench/opt\_*} (KdV, with the nonlinear term $A u u_x$)}\\
Simple         & 20000 &  6 s & $2.4\times10^{-3}$ & --- \\
Natural\_grad  &  3000 & 36 s & $1.5\times10^{-5}$ & --- \\
Kalman         &  1000 &  6 s & $2.6\times10^{-13}$ & $3.8\times10^{-7}$ \\
\bottomrule
\end{tabular}
\end{center}

The KdV residual contains $A\,u\,u_x$ and is not a linear functional
of the carried derivatives. The filter then receives the residual
value directly through \code{kf\_update\_resid} (building the
prediction as $\mathrm{dot}(\text{seed},T)$ via
\code{kf\_update\_obs} would double-count the quadratic term), so
\code{bench/opt\_kalman} also verifies the observation path for
nonlinear operators.

Only the natural gradient carries further keys. The metric is
$G=\frac{1}{N_b}\big[\sum_n \mi{j}_n\mi{j}_n^{\!\top}+\mu\,\mathrm{tr}(\cdot)I\big]$
and the update $\Delta W=-\eta(t)G^{-1}(\nabla/N_b)$. The rows
$\mi{j}_n$ are $\partial u/\partial W$ for \code{DATA} terms (a
Gauss--Newton metric) and $\partial R/\partial W$ for
\code{COLLOCATION} terms; for \code{HOD\_DATA} terms the per-point
loss gradient (an empirical Fisher metric).

\code{LBFGS} is the limited-memory quasi-Newton method: the search
direction comes from the two-loop recursion over the most recent
pairs $(\mi{s},\mi{y})$ (verified by
\code{tools/example\_lbfgs\_check.f90}, $1.3\times10^{-16}$), and the
step from a backtracking line search on the Armijo condition
$L(\mi{w}+t\mi{d})\le L(\mi{w})+c_1 t\,\mi{g}\!\cdot\!\mi{d}$. There
is no learning rate. It is a \textbf{full-batch method}: a pair
$(\mi{s},\mi{y})$ is meaningful only when consecutive gradients
belong to the same function, and a redrawn minibatch swaps the
function at every step; \code{Num\_batch} must therefore equal the
training count (all points minus \code{Num\_validation}), and the
input check stops with the required value on violation.
\code{Lbfgs\_ss /} is self-scaled BFGS (SSBFGS, Oren--Luenberger):
instead of applying the newest pair's scale
$s\!\cdot\!y/y\!\cdot\!y$ once to the initial Hessian, every stored
pair receives its own factor
$\tau_j = s_j\!\cdot\!y_j/(y_j^{\top}H_j y_j)$ at its own level of
the recursion. $H_j$ is the operator built from the older pairs, so
the factors are fixed oldest-first by a nested application of the
same recursion ($O(m^2 n_w)$ per direction, negligible beside one
full-batch gradient). With a single stored pair it coincides exactly
with plain L-BFGS (measured: eight digits over fifty iterations). The
measured verdict is stated honestly: on the problems of this
distribution it is slightly worse than plain L-BFGS
($5.2\times10^{-3}$ against $3.3\times10^{-3}$ at 300 cold EHD
epochs).
\code{Lbfgs\_wolfe c2 /} adds the curvature (second Wolfe) condition
$g(w+t\mi{d})\!\cdot\!\mi{d} \ge c_2\, g(w)\!\cdot\!\mi{d}$
($c_1 < c_2 < 1$) to the line search: a step that lowers the loss
while the slope is still steep is rejected and the step is grown
rather than shrunk (with proper bracketing; each such trial costs a
gradient instead of a cost evaluation). The default 0 reproduces the
pure Armijo search exactly. On the problems of this distribution it
changes nothing --- the first trial $t=1$ already satisfies both
conditions, so the search was not the limiter; the measurement is
recorded with the feature.

\code{Lbfgs\_expand n /} is forward tracking: when the first trial
satisfies Armijo, the step is doubled up to $n$ times and the longest
step that still satisfies the condition and lowers the loss is taken
(default 0 reproduces pure backtracking exactly). Measured on the
EHD cold start it gives a 19\,\% lower loss at equal epochs, but
loses at equal wall time due to the extra cost evaluations. The
memory size is \code{Lbfgs\_m n /} (default 8, range $[1,200]$).
More memory improves the per-epoch quadratic model (the two-loop's
extra cost is $O(n\,n_w)$): on the five-component EHD cold start the
objective at 300 epochs measures $6.3\times10^{-3}$ at $m{=}8$
against $3.3\times10^{-3}$ at $m{=}40$. That case admits
\textbf{one-shot training from random weights}
(\code{zwork/pinn\_ehd/input\_cold.dat}: objective
$3.8\times10^{-4}$ in 900 epochs, about 250 s, $R^2 \ge 0.9975$ in
all components). If no loss-lowering step is found, the memory is
discarded and the run proceeds to the next epoch (\code{Lbfgs\_scan /}
enables extra backtracking before giving up; \code{Lbfgs\_verbose /}
prints the accepted and the recomputed loss at each exit). The
objective the line search itself minimizes is recorded in
\code{lbfgs\_history.dat} (one line = epoch, accepted loss, step,
trial count, acceptance flag); the training history is a different
sum common to every update rule and is left alone.

The natural gradient's damping is chosen by \code{Ngd\_damping}:
\code{TRACE} is $G+\mu\,\mathrm{tr}(G)I$ (relative; its meaning
depends on the weight count), \code{ABS} is $G+\mu I$ (absolute;
independent of it --- the one that works at the network sizes coupled
systems need). Default \code{TRACE}.
\code{Ngd\_dual /} solves the natural-gradient step through the Gram
(dual) matrix: the push-through identity
\[ (\mi{G}+c\mi{I})^{-1}\mi{b}
   = \tfrac{1}{c}\bigl[\mi{b}
   - \tfrac{1}{N}\mi{J}^{\top}\mi{K}^{-1}\mi{J}\mi{b}\bigr],
   \qquad \mi{K}=c\mi{I}+\tfrac{1}{N}\mi{J}\mi{J}^{\top} \]
(with $\mi{K}$ of row dimension) yields \textbf{the identical step}
to the dense $n_w\times n_w$ solve at $O(N^2 n_w + N^3)$ (verified
against the primal route on the EHD case to every printed digit;
140 s/epoch primal against 0.16 s/epoch dual there). The dense 2\,GB
guard does not apply to the dual route.

For a system the metric carries \textbf{one row per residual} (one
per component for a multi-component fit), built with the exact
$\sqrt{w_r}$ factorization of the weighted objective. A single row
per point ($\mathrm{d}(\sum_r w_r R_r^2/2)/\mathrm{d}\mi{w}$) would
collapse the outer product to the rank-one-per-point empirical
Fisher, blind to how the residuals trade against each other --- the
very difficulty of a coupled system. With this and the dual solve,
the natural gradient solves the EHD system from random weights
(minibatch of 40, $0.743\to2.8\times10^{-4}$ in 900 epochs, about
150 s, $R^2\ge0.9986$ in all components;
\code{zwork/pinn\_ehd/input\_cold\_ngd.dat}).

\code{Ngd\_geo h alpha /} is geodesic acceleration
(Transtrum--Sethna): the second directional derivative of every
observable along the step $\delta_1$ is measured by one extra
evaluation, a second solve against \textbf{the same metric} (cheap in
the dual route) gives the correction $\delta_2$, and
$\delta_1+\delta_2/2$ is taken when $\lVert\delta_2\rVert/2 \le
\alpha\lVert\delta_1\rVert$ (an $\eta$-consistent finite difference).
Measured neutral: within trajectory noise on the problems of this
distribution, at about twice the cost per epoch.

\code{Ngd\_trust mu0 /} is a Levenberg--Marquardt rule adapting the
absolute damping from the observed loss change: $\mu\leftarrow\mu/3$
on descent, $\mu\leftarrow 5\mu$ on ascent, bounded to
$[10^{-6},10^{8}]$. Setting it forces \code{Ngd\_damping ABS}.

\begin{longtable}{@{}p{0.36\textwidth}p{0.58\textwidth}@{}}
\toprule
Key & Meaning \\
\midrule
\endhead
\code{Ngd\_schedule\_eta NAME a b /} & the schedule of $\eta(t)$: \code{NONE} (constant $p_1$), \code{SIMPLE} ($p_1/(1+at)^b$), \code{EXP} ($p_1e^{-at}$)\\
\code{Ngd\_schedule\_mu NAME a b /} & the schedule of $\mu(t)$; the same three\\
\code{Ngd\_eta\_bound x /} & lower bound on $\eta$\\
\code{Ngd\_mu\_bound x /} & lower bound on $\mu$\\
\bottomrule
\end{longtable}

%=====================================================================
\section{Output files}

Everything is written to the run directory; build products collect in
\code{build/}.

\subsection{The training history \code{history\_ep<start epoch, 7 digits>.dat}}

One line per validation. The file name contains the starting epoch,
so a restarted run opens a new file rather than appending. The
leading \code{\#} lines are self-describing. Every metric is
normalized \textbf{per point}.

Each metric has two channels: the \textbf{value channel} (\code{\_f})
is the contribution of $N(\mi{x})-y$; the \textbf{derivative/residual
channel} (\code{\_df}) that of the high-order terms and the
collocation residuals. The total columns (3--5) are assembled per
form: \code{DATA} value channel only, \code{COLLOCATION} residual
channel only, \code{HOD\_DATA} $\lambda_0\cdot$value $+$
derivatives.

\begin{longtable}{@{}p{0.14\textwidth}p{0.80\textwidth}@{}}
\toprule
Columns & Content \\
\midrule
\endhead
1 & epoch\\
2 & the patience counter (compared with \code{P\_max})\\
3--5 & total cost: training, validation, all input points\\
6--9 & cost split: training (value, deriv), validation (value, deriv)\\
10--11 & all-points cost split (value, deriv)\\
12--17 & RMSE: training (value, deriv), validation (value, deriv), all points (value, deriv)\\
18--23 & MAE, same layout\\
24-- & only with two or more loss terms: 21 columns per term (the layout of 3--23), in order --- 24--44 for term 1, 45--65 for term 2, and so on\\
\bottomrule
\end{longtable}

A \code{DATA}+\code{COLLOCATION} run, for instance, has 65 columns
per line, with the derivative-channel columns of term 1 (\code{DATA})
zero and the value-channel columns of term 2 (\code{COLLOCATION})
zero.

\subsection{Predictions}

\begin{longtable}{@{}p{0.36\textwidth}p{0.58\textwidth}@{}}
\toprule
File & One line holds \\
\midrule
\endhead
\code{output\_set<4 digits>.dat} & $\mi{x}(1{:}D_0)$, prediction, target. One file per loss term, numbered in \code{Loss\_term} order. For a \code{COLLOCATION} term the target column carries the exact solution (with \code{Exact\_solution 1})\\
\code{output\_hod\_set<4 digits>.dat} & $\mi{x}(1{:}D_0)$, all predicted carried derivatives $T^{\mi\alpha}(1{:}N_\alpha)$, all targets. Written for \code{HOD\_DATA} terms\\
\code{output\_deriv\_set<4 digits>.dat} & $\mi{x}(1{:}D_0)$, $\partial N/\partial x_i\,(i=1..D_0)$. Only with \code{Output\_deriv 1}\\
\bottomrule
\end{longtable}

Since the last two columns are prediction and target, the quality of
a fit is one line of arithmetic away. A coefficient of determination
near one means the model explains the data; near zero it has learned
nothing beyond the target mean. A model outputting a constant still
produces smooth, plausible-looking files, so this check is worth
doing every time.

\subsection{Weights and restarting}
\label{sec:restart}

\begin{longtable}{@{}p{0.36\textwidth}p{0.58\textwidth}@{}}
\toprule
File & Content \\
\midrule
\endhead
\code{nn\_weight.dat} & the best weights (by training cost). Line 1 the best epoch; lines 2--3 reserved fields (activation codes, currently always 0); then the layer count, the widths, and per layer a \code{\#l=} line followed by the rows $W^{(l)}_{j,0:n_{l-1}}$\\
\code{checkpoint\_ep<7 digits>.dat} & the weights every \code{Checkpoint\_cyc}; same format, line 1 its own epoch. Copied over \code{nn\_weight.dat} it becomes a restart point\\
\code{gd\_dw.dat}, \code{gd\_m.dat}, \code{gd\_v.dat}, \code{gd\_r.dat}, \code{gd\_u.dat} & the optimizer state (previous step and the moments); how many exist depends on the method\\
\bottomrule
\end{longtable}

\code{Restart 1} reads the weights from \code{nn\_weight.dat}. If
\code{gd\_dw.dat} is present in the same directory, \textbf{the
optimizer state is restored too}, so Adam resumes without losing its
moments; absent, only the optimizer starts cold. If the epoch of the
weight file and the optimizer log disagree (an old checkpoint copied
in with the optimizer log left behind), the run stops.

\subsection{Reference information}

\begin{longtable}{@{}p{0.36\textwidth}p{0.58\textwidth}@{}}
\toprule
File & Content \\
\midrule
\endhead
\code{hod\_alpha\_order.dat} & the carried multi-index order: slot, $|\mi\alpha|$, $\mi\alpha(1{:}D_0)$ per line. Defines the column order of high-order data files and the slots of the residual echo\\
\code{data\_division.log} & the indices of the points held out for validation\\
\code{hod\_golden.dat} & known values for the self-check; absent, the regression comparison is skipped\\
\bottomrule
\end{longtable}

\section{Verification specification}
\label{sec:verify}

With a single propagation engine, correctness is checked
\textbf{against the definition}: finite-difference verification of
the weight gradient at three levels.

\begin{longtable}{@{}p{0.34\textwidth}p{0.28\textwidth}p{0.30\textwidth}@{}}
\toprule
Level & Object & Measured (max.\ rel.\ error) \\
\midrule
\endhead
\code{tools/dbg\_backward\_fd.f90} (point) & a synthetic dense net ($D_0=2$, $K=3$), arbitrary seed & $2.2\times10^{-10}$ \\
same (batch vs.\ pointwise) & one batch along two paths & $0$ (loop build)\\
same (batch vs.\ FD) & central differences of the batch loss & $1.4\times10^{-11}$ \\
\code{train\_example} test 1 & trained weights, arbitrary seed & $2.9\times10^{-11}$ \\
\code{Hod\_check 1} & in-run, real data and loss & $3.4\times10^{-10}$ \\
\bottomrule
\end{longtable}

Finite differences bound the derivative error but are blind in
principle to defects that leave the converged value intact --- a
write one element past an array, a read of a never-set variable, a
division that is only sometimes by zero. That layer is covered by two
means:

\begin{itemize}[leftmargin=1.4em]
\item \code{make harden}: a build with bounds checking, uninitialized
      values poisoned by a signalling NaN, and IEEE traps (invalid /
      zero / overflow). The shipped cases run under it with zero
      violations.
\item \code{make negtests}: the input-rejection contract and exact
      properties. Malformed or incomplete inputs must stop with a
      diagnostic, and properties that hold exactly must do so (a
      committee of identical members returns exactly zero spread,
      among others).
\end{itemize}

The checks run under both the loop build and the \code{BLAS=1}
build, which compile different kernels (explicit loops against
matrix products). The two compute the same quantities in different
summation orders: the gradients agree to rounding, and four of the
five benchmarks have bit-identical histories, but a long run can
fork once a last-digit difference appears (\code{kdv} is
bit-identical for 60 validation cycles, then forks to a final cost
of the same size). Numbers quoted from runs therefore state which
build produced them. Finite differences are judged relative to the
gradient norm: central differences cannot resolve components orders
of magnitude below the largest, so per-component ratios would
measure the reference's noise rather than the adjoint's error.

The self-check's pass criterion is \textbf{relative to the size of
the gradient itself}. Dividing by $\max(|fd|,1)$ becomes an absolute
tolerance once the gradient drops below one --- on a loss of order
$10^{-3}$ it would test only three digits. And the known-value check
fails on zero comparisons: a \code{hod\_golden.dat} generated with a
different $D_0$, $K$ or closure would have every row skipped as out
of range and could otherwise pass with nothing compared.

The self-check \code{Hod\_check 1} sets deterministic weights and
verifies (a) the carried derivatives against central differences in
$x$ (orders 1--2 only; at $K>2$ the higher slots are covered through
the seed path by (b) and by \code{fdcheck}), (b) the weight gradient
against central differences of the loss, (c) regression agreement
with the known values (\code{hod\_golden.dat}); any miss prints
\code{FAILED}. What is checked is the exact path that ships.

As a regression net, the training-history columns of five benchmarks
(\code{kdv}, \code{zk7}, \code{g7}, \code{opt\_ngd},
\code{opt\_kalman}) are held as frozen references. Thread against
serial accumulation ($3.2\times10^{-15}$), the concurrent training of
several networks, and the agreement of Kalman training with the
ridge-regression solution ($5\times10^{-13}$) are checked by
\code{train\_example} and \code{libverify\_example}.

%=====================================================================
\section{Application to force fields}

\code{tools/example\_hod\_ff.f90} demonstrates a force field trained
on energies, forces and \textbf{analytic Hessians}. Descriptor
pipelines (n2p2 and kin) provide $dG/dR$ but not $d^2G/dR^2$, so they
can train up to forces but not Hessians; this gap is the use that
states the library's position most clearly.

The model is a sum of atomic energies $E=\sum_i N(G^{(i)})$, carrying
the dense set at $D_0=3$, $K=2$ ($N_\alpha=10$). The model Hessian is
assembled by the chain rule,
\[
H_{pq}=\sum_i\Big[\sum_{ab} \frac{\partial^2 N}{\partial G_a\partial G_b}
\frac{\partial G_a^{(i)}}{\partial x_p}\frac{\partial G_b^{(i)}}{\partial x_q}
+\sum_a \frac{\partial N}{\partial G_a}
\frac{\partial^2 G_a^{(i)}}{\partial x_p\partial x_q}\Big],
\]
where $\partial^2N/\partial G_a\partial G_b$ of the first term is the
carried slot $\mi\alpha=\mi e_a+\mi e_b$ and
$\partial N/\partial G_a$ of the second is $\mi\alpha=\mi e_a$. The
loss seeds land on the same slots, so \textbf{Hessian training costs
one ordinary adjoint per atom and configuration}.

\subsection{Verification}

The demonstration contains three hand-written derivative assemblies
(the descriptor derivatives, the model Hessian, the seed of the
Hessian loss). Any of them would train a \emph{plausible} force field
with one wrong coefficient, so each is checked by central differences
before any physics is reported. Measured (relative):

\begin{center}
\begin{tabular}{@{}lr@{}}
\toprule
Check & Max.\ relative difference \\
\midrule
$dG/dx$, $d^2G/dx^2$ vs.\ differences of $G$ & $4.7\times10^{-10}$ \\
model Hessian vs.\ differences of the model energy & $4.4\times10^{-7}$ \\
$dL/dw$ with the Hessian term vs.\ differences of $L$ & $6.1\times10^{-8}$ \\
\bottomrule
\end{tabular}
\end{center}

\subsection{The smoothness requirement (important)}

Because second derivatives are propagated, \textbf{the cutoff
function must be $C^2$ as well}. The cosine cutoff
$f_c=(\cos(\pi r/r_c)+1)/2$ used by the demonstration has
$f_c(r_c)=f_c'(r_c)=0$ but
$f_c''(r_c)=+\tfrac12(\pi/r_c)^2\neq0$: the second derivative jumps
at $r_c$, so a Hessian-trained model is only $C^1$ across the
cutoff.

The demonstration works because no pair of the sampled
configurations approaches $r_c$ --- a property of the sample, not of
the method. The program therefore \textbf{checks that margin rather
than assuming it} (measured 0.24 in equilibrium-spacing units). A
production coupling should use a $C^2$ cutoff (a polynomial, or
$\tanh^3$).

This is the same general rule stated for the activations: \textbf{the
smoothness class of every ingredient must exceed the order of the
derivatives one intends to use}.

\subsection{Result}

Two networks are trained from the same initial weights on the same 60
configurations under the same schedule:

\begin{center}
\begin{tabular}{@{}lrr@{}}
\toprule
Training data & Max.\ rel.\ error, force constants & Max.\ rel.\ error, phonons \\
\midrule
$E$, $F$ (a descriptor pipeline's reach) & 1.85\% & 3.84\% \\
$E$, $F$, $H$ (this library) & 0.55\% & 1.12\% \\
\bottomrule
\end{tabular}
\end{center}

The figure is drawn from the run logs by
\code{bench/post/make\_fig\_forcefield.py}.

\section{Shipped examples and test cases}

\begin{longtable}{@{}p{0.26\textwidth}p{0.68\textwidth}@{}}
\toprule
Where & Content \\
\midrule
\endhead
\code{bench/} & the numerical experiments of the methods paper: KdV, Kawahara, Zakharov--Kuznetsov (orders 3/5/7), Gross--Pitaevskii, the slit, the 10-dimensional EYU case, closure/dense scaling, optimizer comparison, speed comparison\\
\code{zwork/} & the keyword tour: a value fit (\code{x2+1}), high-order fits (\code{hod\_4d\_k3}, \code{hod\_4d\_k7}), collocation (\code{pinn\_kdv}, \code{pinn\_zk7}), coupled flow systems and \code{System} cases (\code{pinn\_lax}, \code{pinn\_kovasznay}, \code{pinn\_cavity}, \code{pinn\_ehd}), a derivative fit (\code{tour\_derivfit}), a Sobolev loss (\code{tour\_sobolev}), learning-rate schedules (\code{tour\_schedule})\\
\code{tools/} & worked examples: embedding (\code{example\_embed}), force fields (\code{example\_mlff}), the training API (\code{example\_train}), verification (\code{example\_libverify}), committees (\code{example\_uq}), custom losses (\code{example\_customloss}), the C examples\\
\code{host/} & use from a host code\\
\bottomrule
\end{longtable}

%=====================================================================
\section{Cost and memory}

With $N_\alpha$ multi-indices and width $n$, one point's forward pass
costs $O(L\,n^2 N_\alpha)$ in the linear part and
$O(L\,n\,(\text{table entries}))$ in the Bell part; the reverse pass
is of the same order. The work space is $O(L\,n\,N_\alpha)$ (the
adjoint plane covers one layer only). $N_\alpha$ is
$\binom{D_0+K}{K}$ for the dense set and falls to the necessary
minimum under a closure.

A value-fit-only run takes the batch path (\code{net\_grad\_batch}),
where the whole minibatch becomes three matrix products per layer
(under \code{USE\_BLAS}).

%=====================================================================
\section{Known limitations}
\label{sec:limits}

\begin{itemize}[leftmargin=1.4em]
\item The hidden activation is $\tanh$ (plus the tabulated
      alternatives of Section~\ref{sec:activation}); carrying another
      activation needs a generator for its derivative table.
\item The read-out is linear only. The scalar forms
      (\code{Residual}, \code{HOD\_DATA}) have one output; the
      \code{System} block (Section~\ref{sec:system}) sets the output
      width to the number of field components and carries the seed as
      a (component $\times$ multi-index) matrix (verified by
      \code{tools/example\_multiout*.f90}).
\item No residual connections, convolutions or attention (fully
      connected only).
\item The primal natural gradient builds a dense metric and is for
      small networks (with a stop on quadratic memory); the dual
      route (\code{Ngd\_dual}) removes that limit for minibatch row
      counts.
\item \code{checkpoint\_ep*.dat} holds weights only; the optimizer
      state lives in \code{gd\_*.dat}. Restarting from an old
      checkpoint therefore requires the matching optimizer log (a
      mismatch is detected and stops the run).
\item The scalar residual's nonlinearity is limited to
      $u\,T^{\mi\alpha}$. The \code{System} block writes products of
      up to three derivative factors (\code{XUX}, \code{DXD},
      \code{TRP}), so $\nabla\rho\!\cdot\!\nabla\phi$ and
      $\rho u u_x$ are expressible; transcendental functions of the
      field, $\sin u$ or $e^u$, are expressible in neither form.
\item Residual coefficients are constants. Space-dependent
      coefficients $c(\mi{x})$ (inhomogeneous media, potential
      terms) are not expressible; since \code{SRC} already reads a
      per-point value, they would extend that mechanism.
\item There is no bare-Laplacian keyword: $\nabla^2u$ is written as
      one \code{LIN} line per dimension (\code{DXLAP} is
      $\partial_x\nabla^{2k}u$ specifically).
\item C wrappers for the committee and the Kalman filter are not
      provided.
\end{itemize}

\end{document}
